% Generated draft: requires mathematical and source review.
\documentclass[11pt]{article}
\usepackage[margin=25mm]{geometry}
\usepackage{fontspec}
\setmainfont{Latin Modern Roman}
\usepackage{amsmath,amssymb,mathtools,unicode-math}
\setmathfont{Latin Modern Math}
\usepackage{enumitem,needspace,xurl,hyperref}
\setlength{\parindent}{0pt}
\setlength{\parskip}{5pt}
\setlength{\emergencystretch}{3em}
\newcommand{\N}{\mathbb{N}}
\newcommand{\Z}{\mathbb{Z}}
\newcommand{\Q}{\mathbb{Q}}
\newcommand{\R}{\mathbb{R}}
\newcommand{\C}{\mathbb{C}}
\newcommand{\F}{\mathbb{F}}
\newcommand{\A}{\mathbb{A}}
\newcommand{\PP}{\mathbb{P}}
\newcommand{\E}{\mathbb{E}}
\newcommand{\eps}{\varepsilon}
\newcommand{\dd}{\,\mathrm{d}}
\begin{document}
\section*{UnsolvedMath ID: 30006348}
\textbf{Source problem code:} catalogue entry
\subsection{Definitions and mathematical statement}
Let \(p\) and \(r\) be primes, and let \(A_p\) be the alternating group on \(p\) letters. For a field \(K\) of characteristic \(r\), a representation is a homomorphism \(\rho\colon A_p\to\operatorname{GL}_K(V)\), where \(V\) is a nonzero \(K\)-vector space. It is faithful if \(\ker\rho=\{1\}\), and irreducible if its only \(A_p\)-invariant \(K\)-linear subspaces are \(0\) and \(V\). For \(g\in A_p\), write \(|g|\) for its order and define its fixed subspace by
\[
V^{\langle g\rangle}=\{v\in V:\rho(g)v=v\}.
\]
The element \(g\) acts without nonzero fixed vectors if \(V^{\langle g\rangle}=\{0\}\).

For primes \(p,r\), the supplied representation question asks whether there exist a field \(K\) of characteristic \(r\), a nonzero \(K\)-vector space \(V\), and a faithful irreducible representation \(\rho\colon A_p\to\operatorname{GL}_K(V)\) such that
\[
\forall g\in A_p,\qquad
\bigl(|g|=p\ \text{or}\ |g|=p-2\bigr)
\Longrightarrow V^{\langle g\rangle}=\{0\}.
\]
No algebraic closure assumption on \(K\), absolute irreducibility requirement, lower bound on \(p\), or primality assumption on \(p-2\) is supplied. The displayed condition quantifies over elements actually having the indicated orders. The record does not settle whether the intended task is to classify all pairs \((p,r)\), to establish existence for some pair, or to answer a question for a restricted family of primes. In particular, the original phrase ``in any characteristic \(r>0\)'' does not unambiguously specify its quantifier over characteristics.

\textbf{Supplied qualifications.} The registry labels this problem open. Its dated curation says that the simultaneous fixed-point-freeness question remains open and that only an obstruction for characteristic \(r>p\) is recorded. No precise statement or hypotheses for that obstruction are supplied here; it is therefore retained as a source qualification, not imposed as an independently verified theorem or as a restriction on the question.

\textbf{Unresolved adjacent source fragment.} The original-statement field continues with a separate ``Problem 2'' about a subgroup \(G\leq\operatorname{Sym}(\Omega)\) having exactly two orbits and possessing a derangement, meaning an element fixing no point of \(\Omega\). Its orbit-size expression is corrupted as \(\texttt{12}\,|\Omega|>1\), and a later expression appears as \(\texttt{21}\,|\Omega|\); these expressions are not silently interpreted as fractions. That fragment attributes the conjecture to Ellis and Harper, reference [1], and reports proofs when \(G\) is simple, nilpotent, has order at most \(1000\), acts primitively on an orbit, or when the corrupted expression \(\texttt{21}\,|\Omega|\) is a prime power. It also attributes to Lee, Popiel and Verret, reference [2], verification for \(|\Omega|\leq60\) and for the corrupted expression \(\texttt{12}\,|\Omega|\) being a product of distinct primes. These are supplied claims about an adjacent problem, with incomplete references and unresolved notation; they do not supply additional hypotheses for the representation question.
\subsection{Short English statement}
The question asks when \(A_p\) has a faithful irreducible representation in positive prime characteristic such that every element of order \(p\) or \(p-2\) fixes only the zero vector. The intended range of primes and quantification over characteristics remain uncertain. The supplied open status and characteristic obstruction require independent review, and the original record also contains a separate, corrupted derangement-problem fragment.
\subsection{Sources}
\begin{itemize}
\item Supplied local registry record, permanent UnsolvedMath ID \(30006348\), title: \emph{Fixed-Point-Free Representations of Alternating Groups}; exact source problem code: \texttt{catalogue entry}. The project identifies \texttt{data/unsolvedmath/problems.json} and \texttt{catalog/problems.json} as its local registries; this draft uses the supplied complete record without independently checking either file. The record lists proposed year \(2025\) and supplies no named proposer.
\item \emph{Computational Group Theory} (2025), Oberwolfach Reports, pp.\ 1389--1450, DOI \(10.4171/owr/2025/27\): \url{https://doi.org/10.4171/owr/2025/27}. This report citation is supplied provenance and has not been independently consulted for this draft.
\item Supplied literature source: \emph{Oberwolfach Report 27/2025}, \url{https://ems.press/content/serial-article-files/51860}.
\item Literature assessment dated exactly \texttt{2026-08-22}: the supplied curation labels the representation question open and says that only the characteristic \(r>p\) obstruction is recorded. This is supplied dated curation, not an independently verified assessment of the literature or mathematical claims.
\item The original-statement field attributes an adjacent derangement conjecture and special cases to Ellis and Harper [1], and further verification to Lee, Popiel and Verret [2]. Full bibliographic entries for [1] and [2] are not supplied, and these attributions concern the adjacent fragment rather than an identified proposer of the representation question.
\end{itemize}
\end{document}
